\documentclass[10pt,a4paper]{amsart}
\usepackage{amsmath,amssymb,amsthm,graphicx}
% Retained preamble for the Stacks Project chapter {CH} (shared project preamble).
% The chapter file pinned-original-source.tex only carries \input{{preamble}} on its first line,
% so no preamble text lives inside the chapter itself.  The Stacks Project uses one single shared
% preamble for every chapter: the snapshot stages it three times (sources/stacks-chow/preamble.tex,
% sources/stacks-etale-cohomology/preamble.tex and focused-seeds/AGP/sources/stacks-snapshot-20261011/
% preamble.tex) and those copies are identical in size.  The block below is copied mechanically,
% line for line, from sources/stacks-chow/preamble.tex lines 14-18 and 154-217: the Xy-pic setup and
% the native theorem declarations together with the mathematical macros.  Nothing was added, removed
% or redefined; the wrapper class and the amsmath/amssymb/amsthm/graphicx base packages are supplied
% by the assembler.
% For commutative diagrams we use Xy-pic
% Chapter: stacks-cohomology.tex -- Cohomology of Sheaves.
\usepackage[all]{xy}

% We use 2cell for 2-commutative diagrams.
\xyoption{2cell}

% Theorem environments.
%
\theoremstyle{plain}
\newtheorem{theorem}[subsection]{Theorem}
\newtheorem{proposition}[subsection]{Proposition}
\newtheorem{lemma}[subsection]{Lemma}

\theoremstyle{definition}
\newtheorem{definition}[subsection]{Definition}
\newtheorem{example}[subsection]{Example}
\newtheorem{exercise}[subsection]{Exercise}
\newtheorem{situation}[subsection]{Situation}

\theoremstyle{remark}
\newtheorem{remark}[subsection]{Remark}
\newtheorem{remarks}[subsection]{Remarks}

\numberwithin{equation}{subsection}

% Macros
%
\def\lim{\mathop{\mathrm{lim}}\nolimits}
\def\colim{\mathop{\mathrm{colim}}\nolimits}
\def\Spec{\mathop{\mathrm{Spec}}}
\def\Hom{\mathop{\mathrm{Hom}}\nolimits}
\def\Ext{\mathop{\mathrm{Ext}}\nolimits}
\def\SheafHom{\mathop{\mathcal{H}\!\mathit{om}}\nolimits}
\def\SheafExt{\mathop{\mathcal{E}\!\mathit{xt}}\nolimits}
\def\Sch{\mathit{Sch}}
\def\Mor{\mathop{\mathrm{Mor}}\nolimits}
\def\Ob{\mathop{\mathrm{Ob}}\nolimits}
\def\Sh{\mathop{\mathit{Sh}}\nolimits}
\def\NL{\mathop{N\!L}\nolimits}
\def\CH{\mathop{\mathrm{CH}}\nolimits}
\def\proetale{{pro\text{-}\acute{e}tale}}
\def\etale{{\acute{e}tale}}
\def\QCoh{\mathit{QCoh}}
\def\Ker{\mathop{\mathrm{Ker}}}
\def\Im{\mathop{\mathrm{Im}}}
\def\Coker{\mathop{\mathrm{Coker}}}
\def\Coim{\mathop{\mathrm{Coim}}}

% Boxtimes
%
\DeclareMathSymbol{\boxtimes}{\mathbin}{AMSa}{"02}

%
% Macros for moduli stacks/spaces
%
\def\QCohstack{\mathcal{QC}\!\mathit{oh}}
\def\Cohstack{\mathcal{C}\!\mathit{oh}}
\def\Spacesstack{\mathcal{S}\!\mathit{paces}}
\def\Quotfunctor{\mathrm{Quot}}
\def\Hilbfunctor{\mathrm{Hilb}}
\def\Curvesstack{\mathcal{C}\!\mathit{urves}}
\def\Polarizedstack{\mathcal{P}\!\mathit{olarized}}
\def\Complexesstack{\mathcal{C}\!\mathit{omplexes}}
% \Pic is the operator that assigns to X its picard group, usage \Pic(X)
% \Picardstack_{X/B} denotes the Picard stack of X over B
% \Picardfunctor_{X/B} denotes the Picard functor of X over B
\def\Pic{\mathop{\mathrm{Pic}}\nolimits}
\def\Picardstack{\mathcal{P}\!\mathit{ic}}
\def\Picardfunctor{\mathrm{Pic}}
\def\Deformationcategory{\mathcal{D}\!\mathit{ef}}

% --- appended 2026-10-12 (Chief): the project own definitions of the reference/slogan/history
% display environments, copied mechanically from the shared Stacks preamble lines 8-12
% (\usepackage{verbatim} plus the three \newenvironment lines, which the chapter preamble keeps
% verbatim from the project preamble). Nothing added, removed or redefined: these environments
% carry no mathematical content and the project itself typesets them as comments. Without these
% three lines a counted statement whose own text uses them fails with Environment slogan undefined.
\usepackage{verbatim}
\newenvironment{reference}{\comment}{\endcomment}
\newenvironment{slogan}{\comment}{\endcomment}
\newenvironment{history}{\comment}{\endcomment}
\usepackage{cleveref}
\crefname{originallemma}{original lemma}{original lemmas}
\Crefname{originallemma}{Original lemma}{Original lemmas}
\makeatletter
\newlabel{lemma-dual-perfect-complex}{{08DQ}{source}{}{sourcepointer-1}{}}
\newlabel{lemma-dual-perfect-complex@cref}{{[originallemma][0][]08DQ}{[1][0][]source}{}{}{}}
\makeatother
\crefname{originallemma}{original lemma}{original lemmas}
\Crefname{originallemma}{Original lemma}{Original lemmas}
\makeatletter
\newlabel{lemma-quasi-separated-cohomology-colimit}{{01FF}{source}{}{sourcepointer-2}{}}
\newlabel{lemma-quasi-separated-cohomology-colimit@cref}{{[originallemma][0][]01FF}{[1][0][]source}{}{}{}}
\makeatother
\crefname{originallemma}{original lemma}{original lemmas}
\Crefname{originallemma}{Original lemma}{Original lemmas}
\makeatletter
\newlabel{lemma-when-jshriek-compact}{{0F5Z}{source}{}{sourcepointer-3}{}}
\newlabel{lemma-when-jshriek-compact@cref}{{[originallemma][0][]0F5Z}{[1][0][]source}{}{}{}}
\makeatother
\crefname{originalreference}{original reference}{original references}
\Crefname{originalreference}{Original reference}{Original references}
\makeatletter
\newlabel{derived-lemma-injective-resolutions-exist}{{013K}{source}{}{sourcepointer-4}{}}
\newlabel{derived-lemma-injective-resolutions-exist@cref}{{[originalreference][0][]013K}{[1][0][]source}{}{}{}}
\makeatother
\crefname{originalreference}{original reference}{original references}
\Crefname{originalreference}{Original reference}{Original references}
\makeatletter
\newlabel{injectives-lemma-derived-products}{{07D9}{source}{}{sourcepointer-5}{}}
\newlabel{injectives-lemma-derived-products@cref}{{[originalreference][0][]07D9}{[1][0][]source}{}{}{}}
\makeatother
\crefname{originalreference}{original reference}{original references}
\Crefname{originalreference}{Original reference}{Original references}
\makeatletter
\newlabel{derived-lemma-leray-acyclicity}{{015E}{source}{}{sourcepointer-6}{}}
\newlabel{derived-lemma-leray-acyclicity@cref}{{[originalreference][0][]015E}{[1][0][]source}{}{}{}}
\makeatother
\title{Perfect complexes are compact on quasi-compact spaces with qc basis}
\author{The Stacks Project Authors}
\date{stacks-project:cohomology.tex:local-snapshot-20261011; GNU FDL 1.2 or later}
\begin{document}
\maketitle
\noindent\textit{Editorial note.} All delivered mathematical text is the author's own, copied verbatim from the pinned Stacks Project chapter file cohomology.tex; no statement or proof step was rewritten, reordered or removed.\par\medskip
\noindent\textit{Editorial note.} The counted span cites the lemmas recorded in assembly.external\_references with canonical tags as navigation pointers.\par\medskip
\noindent\textit{Allowed standard prerequisites (editorial).}\begin{itemize}
\item Duality for perfect complexes. Cohomology of Sheaves, Lemma cohomology-lemma-dual-perfect-complex. Hom(K,M) = H\textasciicircum{}0(X,K\textasciicircum{} or   otimes \textasciicircum{}L M).
\item Cohomology commutes with sums. Cohomology of Sheaves, Lemma cohomology-lemma-quasi-separated-cohomology-colimit. H\textasciicircum{}p(X,-) commutes with direct sums under the qc hypotheses.
\item Compactness of j\_!O\_U. Cohomology of Sheaves, Lemma cohomology-lemma-when-jshriek-compact. Part (b) instance.
\item Leray acyclicity. Derived Categories, Lemma derived-lemma-leray-acyclicity. RGamma(X,M) = Gamma(X, oplus I\_i\textasciicircum{}*).
\end{itemize}\medskip
\ifdefined\hypertarget\hypertarget{sourcepointer-1}{}\fi
\noindent\textit{Source reference.} Original lemma 08DQ: sources/stacks-chow/tags.txt, source lines 10645--10645. Cross-chapter standard prerequisite cited by the author inside the counted span; it is used as a pointer and its own proof is not part of this item.\par\medskip
\ifdefined\hypertarget\hypertarget{sourcepointer-2}{}\fi
\noindent\textit{Source reference.} Original lemma 01FF: sources/stacks-chow/tags.txt, source lines 1868--1868. Cross-chapter standard prerequisite cited by the author inside the counted span; it is used as a pointer and its own proof is not part of this item.\par\medskip
\ifdefined\hypertarget\hypertarget{sourcepointer-3}{}\fi
\noindent\textit{Source reference.} Original lemma 0F5Z: sources/stacks-chow/tags.txt, source lines 18954--18954. Cross-chapter standard prerequisite cited by the author inside the counted span; it is used as a pointer and its own proof is not part of this item.\par\medskip
\ifdefined\hypertarget\hypertarget{sourcepointer-4}{}\fi
\noindent\textit{Source reference.} Original reference 013K: sources/stacks-chow/tags.txt, source lines 1438--1438. Cross-chapter standard prerequisite cited by the author inside the counted span; it is used as a pointer and its own proof is not part of this item.\par\medskip
\ifdefined\hypertarget\hypertarget{sourcepointer-5}{}\fi
\noindent\textit{Source reference.} Original reference 07D9: sources/stacks-chow/tags.txt, source lines 9404--9404. Cross-chapter standard prerequisite cited by the author inside the counted span; it is used as a pointer and its own proof is not part of this item.\par\medskip
\ifdefined\hypertarget\hypertarget{sourcepointer-6}{}\fi
\noindent\textit{Source reference.} Original reference 015E: sources/stacks-chow/tags.txt, source lines 1505--1505. Cross-chapter standard prerequisite cited by the author inside the counted span; it is used as a pointer and its own proof is not part of this item.\par\medskip
\begin{lemma}
\label{lemma-perfect-is-compact}
Let $X$ be a ringed space. Assume that the underlying topological space
of $X$ has the following properties:
\begin{enumerate}
\item $X$ is quasi-compact,
\item there exists a basis of quasi-compact open subsets, and
\item the intersection of any two quasi-compact opens is quasi-compact.
\end{enumerate}
Let $K$ be a perfect object of $D(\mathcal{O}_X)$. Then
\begin{enumerate}
\item[(a)] $K$ is a compact object of $D^+(\mathcal{O}_X)$
in the following sense: if $M = \bigoplus_{i \in I} M_i$ is
bounded below, then $\Hom(K, M) = \bigoplus_{i \in I} \Hom(K, M_i)$.
\item[(b)] If $X$ has finite cohomological dimension, i.e., if there exists
a $d$ such that $H^i(X, \mathcal{F}) = 0$ for $i > d$, then
$K$ is a compact object of $D(\mathcal{O}_X)$.
\end{enumerate}
\end{lemma}

\begin{proof}
Let $K^\vee$ be the dual of $K$, see
Lemma \ref{lemma-dual-perfect-complex}. Then we have
$$
\Hom_{D(\mathcal{O}_X)}(K, M) =
H^0(X, K^\vee \otimes_{\mathcal{O}_X}^\mathbf{L} M)
$$
functorially in $M$ in $D(\mathcal{O}_X)$.
Since $K^\vee \otimes_{\mathcal{O}_X}^\mathbf{L} -$ commutes with
direct sums it suffices
to show that $R\Gamma(X, -)$ commutes with the relevant direct sums.

\medskip\noindent
Proof of (b). Since $R\Gamma(X, K) = R\Hom(\mathcal{O}_X, K)$
and since $H^p(X, -)$ commutes with direct sums by
Lemma \ref{lemma-quasi-separated-cohomology-colimit}
this is a special case of
Lemma \ref{lemma-when-jshriek-compact}

\medskip\noindent
Proof of (a). Let $\mathcal{I}_i$, $i \in I$ be a collection of injective
$\mathcal{O}_X$-modules. By Lemma \ref{lemma-quasi-separated-cohomology-colimit}
we see that
$$
H^p(X, \bigoplus\nolimits_{i \in I} \mathcal{I}_i) =
\bigoplus\nolimits_{i \in I} H^p(X, \mathcal{I}_i) = 0
$$
for all $p$. Now if $M = \bigoplus M_i$ is as in (a), then we
see that there exists an $a \in \mathbf{Z}$ such that $H^n(M_i) = 0$
for $n < a$. Thus we can choose complexes of injective $\mathcal{O}_X$-modules
$\mathcal{I}_i^\bullet$ representing $M_i$
with $\mathcal{I}_i^n = 0$ for $n < a$, see
Derived Categories, Lemma \ref{derived-lemma-injective-resolutions-exist}.
By Injectives, Lemma \ref{injectives-lemma-derived-products}
we see that the direct sum complex $\bigoplus \mathcal{I}_i^\bullet$
represents $M$. By Leray acyclicity
(Derived Categories, Lemma \ref{derived-lemma-leray-acyclicity})
we see that
$$
R\Gamma(X, M) = \Gamma(X, \bigoplus \mathcal{I}_i^\bullet) =
\bigoplus \Gamma(X, \bigoplus \mathcal{I}_i^\bullet) =
\bigoplus R\Gamma(X, M_i)
$$
as desired.
\end{proof}

\end{document}
